\section*{Chapter \ref{ch:mx:benchmark-algorithms} --- Benchmark Algorithms}

\begin{solution}{exo:mx:benchmark-algorithms:1}
$15\,000\times5.50\%=825$ shares in the first bin and $15\,000\times9.38\%=1\,407$ in the last.
\end{solution}

\begin{solution}{exo:mx:benchmark-algorithms:2}
Trading $0.1\,o$ against the others' $o$ makes it $0.1/1.1=9.09\%$ of all volume. To be 10\% it must trade $r/(1-r)=11.1\%$ of the others' volume
(or 10\% of the whole volume, its own included).
\end{solution}

\begin{solution}{exo:mx:benchmark-algorithms:3}
The buyer paid 10.55 against a day's VWAP of 10.75: $-0.20$. Execution: $\tfrac12(10.10-10.00)+\tfrac12(11.00-11.00)=0.05$. Schedule:
$(0.5-0.25)\times10.00+(0.5-0.75)\times11.00=-0.25$: it bought too much in the cheap bin, a lucky forecast error.
\end{solution}

\begin{solution}{exo:mx:benchmark-algorithms:4}
The execution part varies about as much on both kinds of day (standard deviations 0.2 to 0.24 ticks), while the schedule part is four times larger
on high-volume days (2.17 against 0.48 for the static schedule) because their surges reshape the day's volume. The larger the schedule part, the larger
its share of a variance to which the execution part adds a fixed amount.
\end{solution}

\begin{solution}{exo:mx:benchmark-algorithms:5}
Its market-on-close order raised the closing price by 12.6 ticks against the same day without it, so the benchmark rose to meet its fills: the 5.9
ticks ``gained'' are part of that move. Against the arrival price, which it cannot move, it paid 7.3 ticks, and its impact cost was 6.8, the largest.
\end{solution}

\begin{solution}{exo:mx:benchmark-algorithms:6}
Each trades $q=r(o+3q)$, so $q/o=r/(1-3r)=0.2/0.4=50\%$ of the natural volume, and the total volume is $o/(1-3r)=2.5$ times the natural one. At
$r=25\%$, $q/o=0.25/0.25=100\%$ and the volume quadruples; at $r\ge1/3$ the loop has no fixed point and only the orders' sizes stop it.
\end{solution}

\begin{solution}{exo:mx:benchmark-algorithms:7}
382 stock-days (10.1\%) are high-volume. Their tracking errors are 20.0 basis points for the static curve and 18.9 for the level-plus-AR schedule
(17.7 and 17.2 with news on 8\% of days): more news makes busy days more common and worse, and updating helps a little more.
\end{solution}

\begin{solution}{exo:mx:benchmark-algorithms:8}
The VWAP includes the algorithm's own trades, and in the simulated market its buying raised the day's VWAP by 2.0 ticks: its slippage is measured
against a benchmark it moved. Its impact cost against the day without it was 2.7 ticks, and against the arrival price it paid 3.0.
\end{solution}

\begin{solution}{pb:mx:benchmark-algorithms:1}
\textbf{1.} See the definitions; the benchmarks are the day's VWAP, the time-weighted price, the VWAP over the execution window, the arrival price
and the close.
\textbf{2.} It raised the day's VWAP by 2.0 ticks; its impact cost was 2.7 ticks against a reported 0.6.
\textbf{3.} Its market-on-close order raised the close by 12.6 ticks.
\textbf{4.} The impact cost against the same day without the algorithm (common random numbers): it removes the day's own price moves, whose standard
deviation of about six ticks swamps the cost against arrival.
\textbf{5.} Twenty stocks, 250 days, 26 bins; U-shaped profiles, a market factor, a day-level shock, an AR(1) specific deviation, news on 8\% of
stock-days; a high-volume day has more than 1.5 times the stock's average volume of the previous 20 days (331 of 3\,800 test stock-days).
\textbf{6.} $\log v_k=\log(\bar Vp_k)+\ell+e_k$ with $e$ AR(1); after each bin the level is re-estimated by generalised least squares and the last
deviation fades at rate $\rho$.
\textbf{7.} Summation by parts: $\sum(a_k-b_k)w_k=-\sum_{k<B}D_k(w_{k+1}-w_k)$.
\textbf{8.} 17.7 and 13.7 basis points for the static curve, 17.2 and 13.3 for level plus AR.
\textbf{9.} Its additive AR(1), fitted mostly on quiet days, treats a news surge as noise that will revert (20.6 against 18.3 for static PCA-ARMA).
\textbf{10.} $\sum\frac{q_k}{Q}(f_k-w_k)+\sum(\frac{q_k}{Q}-\frac{v_k}{V})w_k$.
\textbf{11.} The first sixteen high-volume stock-days of the test period and sixteen others spread across it; the simulator draws the \omterm{def:mx:why-there-is-a-spread:informed}{noise traders}'
orders in advance from the day's seed, so both schedules meet the same flow.
\textbf{12.} Static: 2.24, 2.17 and 0.21 ticks; dynamic: 1.75, 1.71 and 0.20.
\textbf{13.} 0.66 and 0.73 ticks: the half-spread and the impact that every child \omterm{def:mx:the-limit-order-book:orders}{market order} pays.
\textbf{14.} \emph{Named result}: on the simulator's high-volume days the volume-forecast (schedule) part explains 99.1\% of the static VWAP's
slippage variance and 98.6\% of the \omterm{def:mx:benchmark-algorithms:dvwap}{dynamic VWAP}'s (slippage standard deviations 2.24 and 1.75 ticks); on the other days, 73.9\% and 68.5\%.
\textbf{15.} No: the dynamic schedule did better on 7 of 16 days, the paired difference in squared slippage is 1.5 with a standard error of 1.8, and
one news day carries most of it; the panel's 3\,800 days show a real but small gain (17.7 to 17.2 basis points).
\textbf{16.} $q=r(o+q)$ gives $q/o=r/(1-r)$.
\textbf{17.} The pace rose from 21.1\% to 24.6\% of the natural volume and the impact cost from 6.4 to 10.0 ticks.
\textbf{18.} The order lasted almost twelve minutes instead of five, long enough for the loop to reach its steady state (24.9\%).
\textbf{19.} That its misses on busy days are mostly the volume forecast's, that updating during the day helps a little on average and not reliably
on any given day, and that a better forecast of news days is worth more than a faster update.
\textbf{20.} Of how well the schedule matched the day's volume and of the spread paid, not of what the order cost: the benchmark moves with the
order.
\end{solution}

\begin{solution}{iq:mx:benchmark-algorithms:1}
A VWAP (or a low participation rate) over the day: no \omterm{def:mx:the-almgren-chriss-framework:urgency}{urgency}, a liquid stock, so minimise impact and look like the day's average; implementation
shortfall only with a view that decays.
\iqlookfor{\omterm{def:mx:the-almgren-chriss-framework:urgency}{Urgency} drives the choice; the benchmark follows.}
\end{solution}

\begin{solution}{iq:mx:benchmark-algorithms:2}
Start from the stock's average profile over recent days, add a day-level estimate and an intraday update (a level-plus-AR model, or a common
component across stocks plus a specific ARMA part), with special handling for known events (index rebalances, expiries). Judge it by the VWAP
tracking error it produces, not by the mean squared error of volumes.
\iqlookfor{A profile plus updates; events; a loss that matches the use.}
\end{solution}

\begin{solution}{iq:mx:benchmark-algorithms:3}
Summation by parts: the schedule part equals $-\sum D_k\Delta w_k$, so an error in one bin that is corrected in the next costs only the price move
between them; what costs is being ahead of or behind the volume for long.
\iqlookfor{The identity; the intuition of being ahead or behind.}
\end{solution}

\begin{solution}{iq:mx:benchmark-algorithms:4}
Each algorithm counts the other's prints, so both trade faster than intended ($r/(1-2r)$ of the natural volume) and pay each other's impact. Cap the
rate on the others' volume excluding the other large buyer if it can be identified, lower the rate, or switch to a schedule that does not follow volume.
\iqlookfor{The feedback; a concrete mitigation.}
\end{solution}

\begin{solution}{iq:mx:benchmark-algorithms:5}
The schedule returns target quantities by time (or by volume) from the order, the market state and what has been filled; the child-order layer
tries to meet them and reports fills and market volume back. Keep the schedule deterministic given its inputs so it can be replayed.
\iqlookfor{Targets and feedback; testability.}
\end{solution}

\begin{solution}{iq:mx:benchmark-algorithms:6}
The market-on-close order moved the close: the fills matched the benchmark because the benchmark rose to the fills, and the fund, valued at that
close, bought at a price its own order had inflated; the next day's price may fall back. Look at the imbalance published before the close and the
price after it.
\iqlookfor{Benchmark contamination; the closing auction's imbalance.}
\end{solution}
